\chapter{グリッチライブ壁紙生成システム}

\section{システム概要}

本システムは、グリッチ表現を用いたインタラクティブなライブ壁紙を生成・表示することを目的とする。静止画像を単に背景として表示するのではなく、時間経過および端末の傾きに応じて、パララックス、RGBずれ、ノイズ、ドット崩壊、帯状ノイズなどの表現を変化させることで、動的かつインタラクティブなグリッチ表現を実現する。

本システムでは、壁紙の描画に用いる色情報および各種パラメータを入力として、複数のグリッチレイヤをリアルタイムに生成する。まず、入力された色パレットをもとに背景となる色面および各描画レイヤを構成する。次に、各レイヤに対して矩形ノイズ、ドット崩壊、帯状ノイズ、RGBずれなどの表現を付加する。さらに、スマートフォンのセンサから端末姿勢を取得し、傾き方向および傾き量に応じてレイヤ位置やグリッチ表現の強度を変化させる。

これらの描画処理にはOpenGL ESを用い、複数のレイヤをリアルタイムに合成した結果をAndroidのライブ壁紙として表示する。

また、本システムでは複数の表現パラメータを変更可能とし、利用者の選好に応じて壁紙の外観を調整できる構成とする。
簡易版では、「ザビモノ」、「ザビカラー」、「ゆらモノ」、「ゆらカラー」の4種類のプリセットを用意し、端末操作に対する反応方式と色彩表現の組合せから壁紙を選択できるようにしている。通常版では、これらに加えてRGBずれ量、ノイズサイズ、ちらつき、帯幅、パララックス感度などを個別に調整することが可能である。


\section{システム構成}

本システムは、主に入力データ生成・解析部、パラメータ管理部、センサ入力部、グリッチ生成部、ライブ壁紙描画部、およびユーザインタフェース部から構成される。


\subsection{入力データ生成・解析部}

入力データ生成・解析部では、グリッチ表現の基礎となる色情報を生成・取得する。本研究では、必要に応じてPollinations AIを用いて元画像を生成し、その画像から代表的な色を抽出する。抽出した色情報はパレットとしてJSON形式で保持し、後続の描画処理で利用する。

ただし、本システムでは元画像そのものを直接変形してグリッチ表現を生成することを必須としていない。画像は主として配色を決定するための情報源として利用し、実際のライブ壁紙では、パレットから取得した色情報をもとに色面、ノイズ、矩形、帯状表現などをシェーダ上で生成する。そのため、色情報があらかじめ定義されている場合には、元画像を用いずにグリッチ表現を生成することも可能である。


\subsection{パラメータ管理部}

パラメータ管理部では、グリッチ表現を構成する各種設定値を保持する。主なパラメータとして、使用する色パレット、RGBずれ量、前景ノイズの大きさ、ちらつき量、最前面ノイズの最大サイズ、帯状ノイズの最大幅、パララックス感度、色のばらつき、アクセント色の出現率などがある。

これらの値は基本的に独立して変更可能であり、各パラメータの組合せによって壁紙全体の表現傾向が決定される。また、評価実験などで用いる簡易版では、複数のパラメータをあらかじめまとめたプリセット4種にて管理する。


\subsection{センサ入力部}

センサ入力部では、スマートフォンの姿勢情報を取得し、端末の傾き方向および傾き量を求める。取得した値は、描画レイヤの移動量、RGBずれ量、帯状ノイズのずれ量、ノイズ強度などの制御に利用する。

「ゆら」系の表現では、主として各レイヤの移動量に傾き情報を反映し、複数レイヤ間にパララックスを生じさせる。一方、「ザビ」系では、傾き情報をグリッチ表現の強度へ反映し、端末を動かした際に帯状ノイズ、RGBずれ、ドット崩壊、ちらつきなどを強く発生させる。


\subsection{グリッチ生成部}

グリッチ生成部では、シェーダ内で疑似乱数を生成し、矩形ノイズ、ドット崩壊、帯状ノイズ、RGBずれなどの視覚表現を構成する。

疑似乱数の生成には、画面上の座標、レイヤごとのシード値、時間、端末姿勢に由来する値などを使用する。これらをもとに、各領域におけるノイズの出現、色、位置、形状、欠損量などを決定する。

また、単一種類のノイズを画面全体へ均一に配置するのではなく、異なる大きさや密度を持つ複数のノイズ表現を重ね合わせる。これにより、規則的な図形パターンではなく、複数の表示異常が同時に生じているような複雑なグリッチ表現を生成する。


\subsection{ライブ壁紙描画部}

ライブ壁紙描画部では、Androidのライブ壁紙サービス上でOpenGL ESを用いたリアルタイム描画を行う。

描画は複数のレイヤを奥側から手前側へ順番に合成することで行う。基本構造として5層の描画レイヤを持ち、これに大型矩形ノイズ、最前面のドット崩壊、画面破損オーバーレイ、ザビ系で使用する破片レイヤなどを追加している。

各レイヤにはそれぞれ異なる位置、色、ノイズ量、時間変化を与える。さらに、センサ入力部から得た端末姿勢を反映することで、レイヤ位置および各種グリッチ表現をリアルタイムに変化させる。


\subsection{ユーザインタフェース部}

ユーザインタフェース部では、利用者によるプリセットの選択および各種パラメータの変更を行う。

評価実験等で用いる簡易版では、「ザビモノ」、「ザビカラー」、「ゆらモノ」、「ゆらカラー」の4種類から表現を選択し、ライブ壁紙として設定する構成としている。

一方、通常版では、各プリセットを基礎として、RGBずれ量、ノイズサイズ、ちらつき方式、帯状ノイズの幅、パララックス感度、色のばらつきなどを個別に変更することができる。これにより、利用者が自身の選好に応じてグリッチ表現を調整できる構成としている。

\section{入力データの生成と解析}

本システムでは、グリッチ表現に用いる配色や各種設定を入力データとしてあらかじめ生成し、JSON形式で保持する。入力データを作成する方法の一つとして、生成AIによって作成した画像を利用した。

ただし、本研究では生成画像そのものをライブ壁紙として直接表示するのではない。生成画像から代表色や画像特徴量を抽出し、さらに描画に必要なレイヤ構成やグリッチ設定と組み合わせることで、ライブ壁紙の描画に用いるJSONデータを生成する。

画像の生成および解析は、Android端末上で行うライブ壁紙の描画処理とは分離している。すなわち、画像生成および解析を事前処理として行い、生成されたJSONデータのみをAndroidアプリへ与える構成とした。これにより、端末上で画像解析を繰り返し実行することなく、あらかじめ作成したデータを用いてリアルタイム描画を行うことができる。


\subsection{Pollinations.aiを用いた元画像の生成}

元画像の生成には、画像生成サービスPollinations.aiを利用した。Pollinations.aiでは、生成したい画像の内容をテキストによるプロンプトとして指定し、画像サイズ、生成モデル、seedなどの条件を与えることで画像を生成することができる。

本研究では、特定の物体や風景そのものをライブ壁紙として再現することではなく、グリッチ表現に利用する配色を得ることを主な目的として生成画像を利用した。そのため、画像生成時には、色彩、質感、雰囲気などを表す比較的短いプロンプトを用いた。

例えば、以下のようなURLによって画像を生成した。

\begin{verbatim}
https://image.pollinations.ai/prompt/
misty+cyan+analog
?width=256
&height=256
&model=flux
&seed=123
&nologo=true
\end{verbatim}

この例では、``misty cyan analog''というプロンプトを指定し、256×256画素の画像を生成している。このような語句を用いることで、特定の色調や質感を持つ画像を生成し、後述する解析処理の入力として利用した。

生成画像を利用することで、使用する色を一色ずつ手作業で選択するのではなく、一つの視覚的なテーマを持った画像から、互いに関連した複数の色を取得することができる。例えば、青緑系の雰囲気を持つ画像を生成し、そこから複数の代表色を取得することで、一定の色彩傾向を持ったカラーパレットを作成できる。


\subsection{生成画像の解析}

生成した画像から、ライブ壁紙の入力データとして利用する色情報および画像特徴量を取得する。この解析処理はAndroidアプリとは別に実行し、解析結果を一度データとして保持した後、描画設定と組み合わせてJSON形式へ変換する。

このように画像解析とライブ壁紙描画を分離することで、Androidアプリ側では生成画像そのものを読み込んで解析する必要がなく、あらかじめ作成された描画用データのみを参照できる構成とした。


\subsubsection{色情報の抽出}

生成画像から、画像を特徴付ける複数の代表色を抽出する。本研究で使用したデータでは、5色の代表色について、色を表す16進数形式のカラーコードと、各色の重みを保持する形式とした。

例えば、ある生成画像から得られた色情報は、以下のように表現される。

\begin{verbatim}
"palette": [
  {"hex":"#232d1f", "w":0.0649},
  {"hex":"#baecd6", "w":0.3889},
  {"hex":"#2f5039", "w":0.1754},
  {"hex":"#529795", "w":0.2231},
  {"hex":"#8cd1c3", "w":0.1478}
]
\end{verbatim}

ここで、\texttt{hex}は抽出された色を表し、\texttt{w}は各色の重みを表す。このカラーパレットは、ライブ壁紙における背景色や各描画レイヤ、ノイズなどの色を決定するための基礎情報として利用する。

単一の色のみを用いるのではなく、同一画像から取得した複数の代表色を用いることで、全体として共通した色彩傾向を保ちながら、複数の異なる色をグリッチ表現へ使用することができる。


\subsubsection{その他の画像特徴量}

色情報に加えて、解析結果には画像全体の特徴を数値化した情報も保持する。本研究で作成した解析データには、主として粒状感を表す\texttt{grain}、平均輝度を表す\texttt{luma\_avg}、コントラストを表す\texttt{contrast}を含めた。

例えば、解析結果の一部は以下のように保持される。

\begin{verbatim}
"grain": 0.7306,
"luma_avg": 0.6494,
"contrast": 0.5316
\end{verbatim}

これらは、生成画像の視覚的特徴を後段の処理から参照できるようにするための情報である。色情報と同様に、画像そのものを保持するのではなく、画像の特徴を数値データとして分離して保持することで、描画処理から利用可能な形式とした。

ただし、本研究の中心は画像解析手法そのものではなく、また、現在の描画処理においてこれらすべての特徴量を直接利用しているわけではない。開発過程では、グリッチ表現の外観を確認しながらパラメータを手動で変更する方法も併用している。そのため、これらの画像特徴量は、画像解析結果として保持される情報の一部として位置付ける。


\subsection{描画設定へのパック}

画像解析によって得られた情報は、ライブ壁紙の描画に必要な各種設定と組み合わせ、一つのJSONデータとしてまとめる。本研究では、この処理を画像解析とは別の事前処理として行った。

開発時には、解析結果を含む入力JSONをローカルで動作する処理へ送信し、描画用データを生成した。例えば、以下のようにHTTP POSTリクエストを用いて\texttt{/pack}へ入力データを送信する。

\begin{verbatim}
curl -X POST http://127.0.0.1:8000/pack 
  -H "Content-Type: application/json" 
  -d @analyze_input.json
\end{verbatim}

この処理によって、画像解析によって得られた色情報や特徴量に加え、ライブ壁紙の描画に必要となる背景色、乱数シード、レイヤ構成、グリッチに関する設定などを一つのデータへまとめる。

生成されるJSONデータでは、画像の解析結果を\texttt{analyze}、ライブ壁紙の描画に用いる設定を\texttt{pack}として保持する構成とした。

例えば、\texttt{pack}には以下のような情報を含める。

\begin{itemize}
    \item 背景色
    \item カラーパレット
    \item 乱数生成に使用するseed
    \item 各描画レイヤの識別子
    \item レイヤごとの深度およびパララックス量
    \item レイヤごとの不透明度および更新速度
    \item RGBずれやノイズなどに関するグリッチ設定
    \item 端末の傾きと各グリッチ処理との対応関係
\end{itemize}

例えば、基本となる5層の描画レイヤについては、それぞれに異なる\texttt{depth}、\texttt{parallax}、\texttt{opacity}、\texttt{fps}などの値を保持できる。また、各レイヤで使用する色についても、カラーパレット内の色を異なる組合せで割り当てることができる。

このように、画像解析結果と描画設定をJSON形式の中間データとして保持することで、画像を生成・解析する処理と、Android端末上でグリッチを描画する処理を分離した。また、描画プログラム自体を変更せず、JSON内の値を変更することで異なる配色やレイヤ構成を試すことができる。


\subsection{生成画像の位置づけ}

本研究における生成画像は、最終的なライブ壁紙として直接表示するためのものではなく、主として配色や画像特徴を取得するための入力として利用する。

したがって、生成画像に含まれる物体の輪郭や構図が、そのまま最終的なライブ壁紙へ反映されるわけではない。生成画像から抽出された色などの情報を利用し、ノイズ、矩形、帯状ノイズなどの形状は、後述するシェーダ処理によって新たに生成する。

また、開発初期には、生成画像の解析結果からカラーパレットだけでなく複数のグリッチパラメータまで生成する構成を検討した。一方、実装および表現の調整を進める過程では、各パラメータがグリッチ表現へ与える影響を個別に確認する必要が生じた。そのため、生成された設定をそのまま利用するだけではなく、JSON内の値を直接変更し、ノイズ量、RGBずれ量、パララックス量などを調整しながら描画を確認する方法も併用した。

このため、現行システムにおいて、生成AIによる元画像の作成はライブ壁紙を動作させるための必須処理ではない。必要なカラーパレットおよび描画パラメータがあらかじめ用意されている場合には、生成画像を介さずにJSONデータを作成し、グリッチ表現を生成することも可能である。

したがって、本研究では、生成画像をグリッチ表現そのものではなく、配色や初期設定を作成するための入力手段の一つとして位置付ける。


\section{解析結果に基づくグリッチ生成}

Androidアプリでは、前節までの事前処理によって作成されたJSONデータを入力としてグリッチ表現を生成する。

Android側では、元となった生成画像を再度読み込んだり、画像から色情報や特徴量を抽出したりする処理は行わない。事前に作成されたJSONから、カラーパレット、レイヤ構成、各種描画パラメータなどを取得し、OpenGL ESによるリアルタイム描画へ反映する。

本システムでは、元画像の画素を直接移動、欠損、変形させることでグリッチを生成するのではなく、JSONから取得した色情報や設定値をもとに、シェーダ上でノイズ、矩形、帯状ノイズなどを手続き的に生成する。

この構成により、同一のカラーパレットを使用しながら、ノイズの大きさ、密度、RGBずれ量、パララックス量などのみを変更するといった調整が可能となる。また、画像生成および解析処理と描画処理が分離されているため、Android側の描画処理を変更することなく、異なるJSONデータへ切り替えて表現を変更できる。


\subsection{JSONデータの読み込み}

描画に使用するJSONファイルは、Androidアプリの\texttt{assets}ディレクトリに配置する。

アプリでは、対象となるJSONファイルの内容を文字列として読み込み、Gsonを用いて\texttt{Pack}型のデータへ変換する。描画部では、変換されたデータから、カラーパレット、レイヤ情報、および描画に必要な各種設定値を参照する。

複数のJSONファイルを用意することで、描画処理そのものを変更することなく、異なる配色やレイヤ構成を持つデータへ切り替えることができる。また、開発中の表現検証では、JSON内の値を直接変更することで、ノイズ量、RGBずれ量、パララックス量などを調整し、各設定値が最終的な外観へ与える影響を確認した。

JSONファイルの読み込みまたはデータへの変換に失敗した場合には、エラー内容をログへ出力し、読み込み結果として\texttt{null}を返す構成としている。

\subsection{色情報のグリッチ表現への反映}

\subsection{各グリッチパラメータへの反映}

\section{Androidライブ壁紙への描画}

\subsection{OpenGL ESによる描画}

\subsection{ライブ壁紙としての実装}


\subsection{描画更新}

\section{端末操作に応じた動的表現}

\subsection{センサ情報の取得}

\subsection{傾きに応じた視差効果}

\subsection{グリッチ表現への反映}

\section{ユーザインタフェース}
% 壁紙選択
% パラメータ調整
% プレビュー等

\section{動作環境}
